% =====================================================================
%  R. Nielsen, PROPAEDEUTICS AND LOGIC.  COURSE FOR THE ACADEMIC YEAR 1868-69.
%  Copenhagen: Gyldendal (F. Hegel); printed by J. H. Schultz, 1868.  (Propædeutik og Logik. Cursus for
%  Universitetsaaret 1868-69.)
%  ENGLISH TRANSLATION OF THE NEW PAGES ONLY, from transcription.tex (Danish), 2026-10-13, made in the main session
%  from the collated and fully page-checked transcription; translation.tex is hand-made, not generated: edit it
%  directly.  The rest of the volume is the earlier books in the same setting (see the Editorial note): their English
%  is in the translations of those books.  Where these pages are the same type as an earlier book (p. 81 to the head
%  »B.«; p. 89 from »1) Ordet og Gud«; p. 263; the first lines of p. 264), that book's English is used, so that the
%  same Danish reads the same in English.
%  Conventions: TRANSLATION-PLAYBOOK.md; terms in ebook-method/propaedeutik-1868/GLOSSARY.md, i.e. the Grundideernes
%  Logik glossary (Viden = knowledge, Magt = power, Slutning = syllogism, Virkelighed = actuality, Tilværen =
%  determinate being, Magtsprog = power-language, Phantasie = imagination), with Ideelære = doctrine of the idea,
%  Ideernes Logik = logic of ideas, Oversigt = Synopsis (Hans, 2026-10-13).  Page numbers of the print in the
%  margin; the Synopsis is separately paginated [1]-6.
% =====================================================================
\documentclass[12pt,a4paper,openany]{book}
\usepackage[utf8]{inputenc}
\usepackage[T1]{fontenc}
\usepackage{libertinus}
\usepackage{libertinust1math}
\usepackage{textalpha}
\usepackage[english]{babel}
\DeclareUnicodeCharacter{221E}{\ensuremath{\infty}}  % ∞ (Synopsis p. 1)
\usepackage[protrusion=true,expansion=true]{microtype}
\usepackage{setspace}
\usepackage[margin=1.3in]{geometry}
\usepackage{fancyhdr}
\usepackage{hyperref}
\hypersetup{pdftitle={Propaedeutics and Logic. Course for the Academic Year 1868--69 (the new pages)}, pdfauthor={Rasmus Nielsen}, hidelinks}
\setlength{\headheight}{15pt}
\pagestyle{fancy}
\fancyhf{}
\fancyhead[LE,RO]{\thepage}
\fancyhead[C]{\textit{R. Nielsen: Propaedeutics and Logic (1868)}}
\renewcommand{\thefootnote}{\ifcase\value{footnote}\or *)\or **)\or ***)\else
  \arabic{footnote})\fi}
\setstretch{1.3}
\usepackage{marginnote}
\newcommand{\opage}[1]{\marginnote{\footnotesize\textit{[#1]}}}
\newcommand{\editorial}[1]{\par\medskip\begin{center}\small\textit{[#1]}\end{center}\medskip\par}
\newcommand{\contnote}[1]{\begingroup\renewcommand{\thefootnote}{}\footnotetext{\textit{[continued from p.~262:]} #1}\endgroup}

\begin{document}

\begin{titlepage}
\begin{center}
\vspace*{4em}
{\Huge\textbf{Propaedeutics and Logic.}}\\[2em]
{\Large\textbf{Course for the Academic Year 1868--69.}}\\[3em]
{\small by}\\[1em]
{\Large\textbf{R. Nielsen.}}\\[8em]
{\small\textbf{Copenhagen.}}\\[0.5em]
{\small Published by the Gyldendal Bookshop (F. Hegel).}\\[0.5em]
{\footnotesize Printed by J. H. Schultz.}\\[0.5em]
1868.\\[3em]
{\small Translated from the Danish}
\end{center}
\end{titlepage}

\chapter*{Editorial note}
\addcontentsline{toc}{chapter}{Editorial note}
{\small This volume, Nielsen's course book for 1868--69, was made up for the most part of sheets of his earlier books,
reissued in the same setting.  Only the pages composed for 1868 are translated here; the others are in the
translations of the earlier books:
\begin{center}
\begin{tabular}{ll}
1868 & is \\ \hline
Synopsis pp.~[1]--6 & new (below) \\
pp.~1--80, p.~81 to the head »B.« & \emph{Lectures on ``Philosophical Propaedeutics'' 1860--61}, pp.~1--81 \\
p.~81 from »B.«, pp.~82--88, p.~89 to »1) The Word and God« & new (below; pp.~81 and 89 given whole) \\
pp.~89 (from »1) The Word and God«)--262 & \emph{The Logic of the Fundamental Ideas II} (1866), the same pages \\
pp.~263--264 & the seam into section C (below; p.~264 is new) \\
pp.~265--456 & \emph{The Logic of the Fundamental Ideas I} (1864), the same pages \\
\end{tabular}
\end{center}
Page numbers in the margin are those of the 1868 print.\par}

% --- p. 1 ---
\opage{1}\begin{center}{\large Synopsis.}\end{center}

We open a view into the peculiarity of philosophical knowledge by beginning to apprehend philosophy as the doctrine
of contradiction, as the doctrine of possibility and as dialectic. To characterize the relation of dialectic to the
exact analyses, a survey is added of the relation between the mathematical and the philosophical treatment of the
problem of infinity. From this survey the following main points are brought out.

a) Every finite magnitude can be divided into infinitely many infinitely small parts, and the sum is given by
0.\,∞. Now since the one finite magnitude is nonetheless different from the other, it follows from this that the one
0.\,∞ must be able to be different from the other, that in other words 0.\,∞ is an indeterminate magnitude.

b) The indeterminate magnitude 0.\,∞ requires a determination of the infinitely small and the infinitely great.
This determination rests on the difference between fixed and movable limits. The fixed limit cannot be exceeded; the
movable one can, as regards diminution as well as augmentation, be exceeded into the infinite.

c) Since a quotient of two variable magnitudes remains unchanged when numerator and denominator are either divided or
multiplied by the same magnitude, it is seen from this that quotients of infinite (small or great) magnitudes must be
able to be just as determinate as those of finite ones.

On these presuppositions rests the analysis of the infinite, whose yield for philosophy is that the contradictions
between the finite and the infinite, which are otherwise treated in all too abstract generality, find a punctual
dialectical solution on an exact foundation.

% --- p. 2 ---
\opage{2}\begin{center}\textbf{Propaedeutics and Logic.}\end{center}

The universal foundation of philosophy is logic. This science can be treated partly as formal logic, partly as
ontology, metaphysics or the logic of ideas.

\begin{center}\textbf{Formal Logic}\\
\textbf{in}\\
\textbf{Outline.}\end{center}

Formal logic, or the doctrine of the abstractly necessary forms of thought, divides into the doctrine of concept,
judgment and syllogism.

\begin{center}\textbf{Concept.}\end{center}

Thinking, which is already present in sense-perception, lifts itself out by dissolving the immediate sense-images,
forms by comparison and abstracting a general image, holds it fast in the word, develops its necessarily coherent
essential marks and raises it to a concept (\emph{concept-analysis}; definition).

The sum of the marks indispensable to a concept is called the concept's \emph{content}; the number of the objects to
which it applies is called its \emph{extension} (relation between extension and content; superordinate and
subordinate concepts; concepts of genus and of species).

\begin{center}\textbf{Judgment.}\end{center}

The concept finds its pure logical expression only in language; words are joined into propositions, and in
propositions judgments are uttered. The proposition which by its form expresses a concept is a logical judgment. To
express the concept, the judgment must formally state: the \emph{subject}, the representation that is to be
determined; the \emph{predicate}, the determining representation; and the \emph{copula}, the determining connecting
word.

From the connection of the subject with the predicate, determined by the copula, it follows that the judgment can be
divided into a plurality of logically different % [lo-/gisk at the page turn]
% --- p. 3 ---
\opage{3}judgments, and judgments are \emph{divided}: according to \emph{quality} into affirmative (a is b), negative
(a is not b) and infinite (conscience is no square); according to \emph{quantity} into singular (this a is b),
particular (some a are b) and general (all a are b); according to \emph{relation} into categorical (a is, according
to its category, b), hypothetical (if a is b, then c is d) and disjunctive (a is either b or c or d); according to
\emph{modality} into assertoric (a is actually b), problematic (a is possibly b) and apodictic (a is necessarily b).

If the predicate follows from the subject with general-logical necessity, the judgment is called \emph{analytic}; if
the predicate can be added to the subject only by a particular experience, the judgment is called synthetic
(synthetic judgments a posteriori and a priori).

\begin{center}\textbf{Syllogism.}\end{center}

The judgment must be grounded, its validity seen; the act of thought by which the validity of a judgment is derived
from one or more other judgments is a logical syllogism. If the syllogism has only one premise, it is called
\emph{immediate}; if it has several premises, it is called \emph{mediate} (the requisites of the syllogism).

The \emph{mediate} syllogisms are \emph{divided} into: the \emph{categorical}, which are conditioned by the mark of
the subordinate concept (nota notæ) being subsumed under the superordinate (syllogistic figures); the
\emph{hypothetical}, which rest on the dependence of the conditioned on its condition (modus ponens and modus
tollens); and the \emph{disjunctive}, which are grounded in the members of a whole being so related to one another
that when one is expressly posited, the rest are excluded (modus ponendo tollens), and when the rest are excluded, the
one remaining must be posited (modus tollendo ponens).

By treating the doctrine of proof and method, formal logic resolves itself into a general doctrine of cognition and
of science.

% --- p. 4 ---
\opage{4}\begin{center}\textbf{Main Propositions}\\
\textbf{of the Logic of the Fundamental Ideas.}\end{center}

a) The validity that can be ascribed to philosophy rests essentially on the validity that can be ascribed to the word.
The human word expresses only a powerless knowledge; therefore concept and object fall outside each other. It becomes
in this respect the task of philosophy to show the relation between the derived and the original, between the human
and the divine word.

b) The divine word is: α) in the form of \emph{knowledge}, the inner pictorial that bears representations and
thoughts, concepts and judgments; β) in the form of \emph{power}, the outward imprinting, the factual utterance of the
unity of the concept and the object; γ) in the \emph{reciprocal determination} of knowledge and power, the
teleological unity-in-difference of thinking and acting, judging and doing, comprehending and creating.

c) In the form of \emph{knowledge} the relation between judgment and object is apprehended theoretically. The unity
of the judgment with the object is explained: α) by a \emph{parallelism} (Plato, Cartesius), such that things are
apprehended as sensuous copies of intellectual archetypes; β) by \emph{impression} and \emph{abstraction} (Hume), such
that concepts become intellectual copies of the objects given in experience; γ) by a \emph{subjective determination}
(Kant), such that the subjective forms of cognition, which are affected by an unknown objective being (das Ding an
sich), are assumed to determine the phenomena of experience.

d) In the form of \emph{power} the relation between judgment and object is to be apprehended practically. The unity
of the judgment with the object then requires: α) a \emph{preformation} (ideal praxis), by which the real concept,
the concept that is to be realized, determines the object beforehand with a view to the conditions under which the
realization is to take place; β) a \emph{realization} (ideal-real praxis), by which the universal in the real concept
is translated point by point into its corresponding individual, so that the object % [Gjen-/standen at the page turn]
% --- p. 5 ---
\opage{5}which is brought forth is to be regarded as the concept's practical definition; finally γ) an explication
(return from praxis to theory), by which the one who judges and brings forth shows himself to recognize the concept of
what is brought forth in the object brought forth (with a view to creation).

e) The teleological unity of the concept in the form of \emph{knowledge} with the concept in the form of \emph{power}
rests on interpreting functions, corresponding to reciprocal determinations of \emph{being and signifying}. The
dialectic of means and ends which develops here lets us see that actuality expresses concepts by a categorical
power-language.

f) How the subjective in the \emph{logic of ideas} is related to the objective can be further elucidated by the
development of the concepts possibility, actuality and necessity, which corresponds to the division of judgments
according to modality in formal logic.

\begin{center}\textbf{Propaedeutics and the Doctrine of the Idea.}\end{center}

The idea, the absolute unity of concept and content --- the principle developed in a totality of means and ends, the
identity of the subjective and the objective --- can be said to be both one with and different from actuality. In
opposition to the phenomenon, the outer, the idea is the essence, the inner; and yet the idea is raised above the
separation of essence and phenomenon, of inner and outer; it is, in its opposition to the phenomenon, the essential
unity of the phenomenon and the essence, and in its opposition to the outer of actuality, the actual unity of outer
and inner.

Since the idea thus permeates and determines all actuality, it follows from this that the doctrine of the idea cannot
gain reality except on a comprehensive foundation of logically introductory analyses of actuality. We have examples
of this in the logic of objectivity, in the doctrine of the functions of the ground, of the cause and of the
ground-cause.

\begin{center}\textbf{Main Propositions Belonging to the Doctrine of the Idea.}\end{center}

1) The ground-cause, the principle, or, as we might also call it, the idea-cause, develops its content in totalities,
and \textbf{the development of the totalities} passes through extremes and % [Slut-/ninger at the page turn]
% --- p. 6 ---
\opage{6}syllogisms. Characteristic in this respect is the difference between the \emph{mechanical}, the
\emph{chemical-physical} and the \emph{organic} totality, whose logical development of form through corresponding
functions states the relation between the \emph{cosmological}, the \emph{geological} and the \emph{biological} idea
and forms the foundation of the philosophy of nature.

2) The opposition between nature and spirit is determined in the doctrine of the idea thus: nature, essentially seen,
is the idea in the form of necessity, spirit the idea in the form of freedom. Nature and spirit are one in the idea;
but this fundamental unity must nevertheless be apprehended in such a way that both, by virtue of the idea, are at the
same time originally opposed to each other. Spirit can as little be derived from nature as nature from spirit. The
unity-in-opposition is a syllogism-unity, which is known by the inner infinity that the beginning itself is a
syllogism just as much as the result.

3) The division of \emph{the philosophy of spirit} (after Hegel) into the three sections: a) spirit in its concept
(anthropology, psychology); b) \emph{practical} spirit (philosophy of right, philosophy of history) and c)
\emph{absolute} spirit (aesthetics, philosophy of religion) has, regarded as a schema, its advantages, but must
nevertheless, with respect to principle and method, be subjected to a thorough critique.

\editorial{End of the Synopsis.  The text pp.~1--80, and p.~81 to the head »B.«, is \emph{Lectures on
``Philosophical Propaedeutics'' from the Academic Year 1860--61} (1862), pp.~1--81, in the same setting: see that
translation.}

% --- p. 81 ---
\opage{81}while yet the abstract quantum which furnishes quantifying's immediate presupposition owes its existence
precisely to the quantifying, and thereby comes to presuppose that of which it was itself to be the presupposition
--- it is seen from this that quantity as an abstract concept can first acquire its true significance in and through
an existing thing whose quantity it is.

Of pure mathematics it is said that it finds ``application'' to actual objects (in mechanics, astronomy, optics and
the like). Essentially understood, such an ``application'' presupposes that the mathematical is originally realized
in the objects to which mathematics is applied. An application of mathematics to actual objects would be, if not a
sheer impossibility, at any rate a subjective arbitrariness, if the mathematical did not itself belong with inner
objective necessity to the essence of actuality; but as a moment in the essence of actuality the mathematical
belongs also to philosophy.

\begin{center}\rule{4em}{0.4pt}\\[0.5em]
B.\\[0.3em]
\textbf{Propaedeutics and the Logic of Ideas.}\\[0.5em]
\rule{4em}{0.4pt}\end{center}

In order to recognize the significance which the word as word must exert on the fundamental representations of the
eternal, on the first natural formation of the ideas in consciousness, we note especially the becoming of the
categories in mythology and theosophy.

The distinguishing mark of the mythical is, in principle, by no means unclarity about the outward finite; it is not
the misty images of sensuousness but the lightning-flash of the idea, not the night of the understanding but the
dawn of reason, that lends the content of consciousness its wondrous mythical colours. How peculiar the mixture of
colours is strikes the eye when one would rightly characterize mythology and the mythical standpoint. ``For
mythology,''
% --- p. 82 ---
\opage{82}says Heiberg, ``there is no other, all-embracing point of view than the poetic. Mythology is poetry, and
purely original folk-poetry.'' Later on this folk-poetry is then also apprehended as a kind of speculative
philosophy. ``Every mythology,'' it says, ``begins with cosmogony; for the first endeavour of philosophizing mankind
is to explain the wonder of existence. The first question is: how has the world come to be? What was there before
all this came to be? Or what was there when there was as yet nothing? The understanding can so little answer these
questions that it must even declare them foolish and self-contradictory. But imagination soars boldly over the
beginning of things, without fearing to lose itself in nothing. It peoples the emptiness, it lets light arise out of
darkness, existence out of non-existence''\footnote[1]{J. L. Heibergs nordiske Mythologie. Paa Dansk udgivet af
Chr. Winther. Kbhvn. 1864. p.~4 and p.~24.} etc.

That mythology is at one and the same time both poetry and philosophy will then mean, in other words, that at bottom
it is neither poetry nor philosophy: not poetry, because it is philosophy; not philosophy, because it is poetry. The
closer account is seen from the functions. The poetic is not here apprehended as poetry. Poetry presupposes a
separation of illusion and actuality bearing on the absolute content, a separation that expressly sets illusion above
actuality because it sets the idea above finitude; but with so reflected a separation myth as myth cannot possibly
subsist: the mythical figures hold themselves by the magic of the idea at the stage of a higher actuality; at this
stage there is not even an inkling of illusion. Just as little is the philosophical in the formation of myths
apprehended as philosophy. What kind of philosophy would that be, which began with an express rift between
understanding and imagination, and which on account of this rift gave up the understanding and let imagination
% --- p. 83 ---
\opage{83}rule? Instead of beginning with a dualism of imagination and understanding, the formation of myths begins,
on the contrary, with an imagination impregnated by understanding and reason; one can say with full right of the
mythical idea-movement that thought imagines because \emph{imagination thinks.}

When, then, the imagination impregnated by understanding and reason ``soars boldly over the beginning of things,
without fearing to lose itself in nothing,'' what steers its flight, what is it that keeps it from plunging into the
abyss of emptiness? Essentially it is the \emph{idea}; immediately seen, it is the \emph{word.} The abyss of the
beginning is equally deep whether we begin with Chaos or with \emph{Ginnungagap}; it is infinitely deep in both: but
in the infinite deep the idea is wedded to the word.

If we begin with Chaos, it is of course not the empty name on which weight is laid here, for the names change. But
the synonymous names all sound upon the same fundamental word. That Chaos as fundamental word is a revelation of the
idea, posited and moved by the power of the idea, can easily be shown. Chaos in the mythical sense is neither an
abstract logical concept nor an expression for a condition of physical confusion, but a totality-unity held fast in
the word, a subjective-objective first, which in pictorial and yet non-pictorial form makes intuitable what cannot be
intuited, namely the origin of things, the idea of all actual beginning or the idea of the actual beginning.

That Chaos as a concept is not in the form of the concept but in the form of a non-pictorial picture follows
expressly from the unity of intuition with the word. What is there in Chaos? Nothing, and yet everything! This
self-contradiction of thought, with the consequences springing from it, is not apprehended in the form of a thought
but is seen in a picture, a mighty thought-picture, formed by the word. Since thought, with the help of the word,
imagines over nothing and everything, imagination thinks essentially of everything and yet intuits in Chaos nothing,
the empty; it thinks essentially of nothing, and yet intuits in Chaos everything, namely the seeds of everything
(non bene junctarum discordia semina rerum). The content,
% --- p. 84 ---
\opage{84}the contentless content in Chaos, is essentially subjective by means of the word: a product of imagining
thought, of thinking imagination; but it is at the same time intuited objectively, strikes the intuiting
consciousness with the power of objectivity in the word, and takes on the figure of physical immediate reality.

But that Chaos in the objective sense nevertheless has no physical immediate reality is precisely a consequence of
the non-pictorial in the picture, of the mythical unity of intuition with the word. In an objective physical
significance Chaos would of course have to be the actual dissolution in which the whole of existence must find itself
when all forms were destroyed, and the swarm of the atoms in lawless confusion. For an understanding, critical
consciousness of such a Chaos the word has lost its magical force. The word Chaos is in this case only a name; Chaos,
on the other hand, as an outward actual condition is the thing, and the thing, not the word, is that on which it
depends. Separated from the imagination that rests on the word, the understanding will then on its own account
investigate the actual Chaos, the outward actual beginning of things. But this undertaking of the understanding will
totally miscarry. As an actual condition Chaos is altogether non-actual; there is actually no existence that can
begin with an actual Chaos: what the understanding is to get out of such a Chaos it must in actuality itself smuggle
into it: it is an idle game the understanding plays with itself and with the actual.

In the mythical Chaos objectivity is subjective; the objective content is translated by the word into subjective
form. How different the mythical Chaos is from the outward physical one can also be recognized by the figures that
come forth from Chaos. ``First is Chaos, then the broad Earth, dark Tartarus --- and Eros, the fairest among the
immortal gods.'' He who sees the old gods step forth from Chaos cannot be in doubt what kind of being Chaos is.
Inasmuch as the beginning is made here with Chaos, a run is in a way taken here at an objective presentation of the
coming-to-be of the world, of the formation of things, of the finite realities;
% --- p. 85 ---
\opage{85}but the view turns round. Cosmogony turns over and becomes theogony. Chaos becomes a divinity, the divinity
from whose womb the particular gods are born. This return of the cosmogonic into the theogonic is in a high degree
characteristic of the formation of myths. For the world to come to be, the gods must first come to be; for the finite
realities in actuality to be able to come to be, the real ideas must first come to be: the finite becoming of the real
ideas reflects in itself the eternal being of the fundamental ideas.

Chaos is thus in its ideality eternal, in its immediate reality temporal, divinity in being, ground of nature in
becoming, totality in original unity, ground of the totalities in original plurality; in Chaos the relative is
absolute, the absolute relative; in the Chaos-word motion is rest, and rest again motion; fullness emptiness, and
emptiness fullness; the inward an outward, and this outward again an inward: in Chaos the idea is word, and the word
idea.

It is then not a contingency but an essential necessity that the oldest gods go forth from Chaos. In logical form the
necessity is recognizable in this, that plurality comes forth from unity, since unity itself is originally plurality;
that the particular in the totality comes forth from the universal, since the universal is originally a totality of
particulars. But in the logical form the content lacks immediate reality; reality, mythical reality, comes with the
word, for the word determined by the idea determines the idea.

Such figures as Gaia, Tartarus, Erebus, Nyx and Eros are in their originality more than empty personifications.
``Figures of the figureless,'' ``elementary forms of the first One, in themselves formless,'' such fasting
abstractions fall just as far below the idea as below actuality. No, by the magic of the idea in the mythical word the
oldest gods are, despite what is formless in the form, typical figures, full of reality, each divinity a totality in
itself,
% --- p. 86 ---
\opage{86}a being of its own, a definite physiognomy. It is the idea that in particular totality-unities permeates,
moves and determines the word; but it is in the immediate sense the word, the particular fundamental word, that in its
unity takes up, determines and gathers the manifold of ideal and real, non-pictorial and pictorial moments that make
up the absolute content of the idea.

The goddess Gaia, who takes her name from the earth, takes up by this name all the reality of the earth into herself,
becomes by this name a divinity with immense capacities, mighty forces and great possibilities; but that the goddess
has taken her name from the actual earth says at the same time that it is the actual earth that has taken its name
from the goddess. The earth-goddess and the actual earth have not merely one name, but they are one being, designated
by one word: they are one in the word, because they are one in the idea. That the actual earth is a goddess shows us
in a really striking way what it means that the objective, real content of the idea is --- by means of the word ---
translated into subjective, ideal form; that the goddess with all the traits of elementary subjectivity is the earth
itself, the actual earth, makes intuitable on the other side what the word is capable of when it is a matter of
translating a subjective content --- for the goddess's being is in its ideality a subjective being --- into
immediate, objective form.

The colossal boldness with which the thinking imagination, impregnated by the idea, defies probability at the sight
of active gods is an eloquent testimony to the power of the word in the objectified events. The age of the gods, the
distant past, is the golden time of the ideas, of the ideal causes; the present with its facts is the time of
actuality, of the conditions and of the material causes: can the two times well become one time? In the mirror of the
word, i.e.\ of narrative, of memory, the past is indeed reflected in the present as the great fundamental ideas are
in the actual results of the development; but actuality with its causes in things reacts with the power of the
present against those
% --- p. 87 ---
\opage{87}ideal presuppositions; the consequence is a rift between the thing and the word, a rift between the
objective and the subjective, a rift between content and form, imagination and thought, a rift that awakens
criticism, sublates the standpoint and sets the idea-forming functions new and higher tasks.

So much for the mythical. For the elucidation of the theosophical form we cast a glance at Philo.

``When one,'' says a recent author, ``occupies oneself with Jewish Alexandrianism, one is involuntarily referred to
Hebraism and Platonism, as those under whose influence it has indisputably formed itself. In this all interpreters
are agreed; the difference first appears when one is to make clear to oneself what character of spirit it was that
expressed itself partly in this receptivity itself, partly in the freer productivity which this called into life. In
Gfrörer and Dähne we now find the latest and most marked representatives of the view of it as reflecting; for
although they polemicize sharply against each other and end in practically opposite results, they nevertheless
present the standpoint on which they think of Philo, as well as his predecessors in Alexandrian theosophy, as
standing, as essentially the same. --- As the spiritual forces whose collision brings movement into Philo's spirit
they regard, from the one side, the definite appropriation of the Platonic opposition between the divine and matter,
and of the thought necessarily following from it, that God could not stand, not even in the remotest connection with
finitude, so that even all predicates taken from it were to be denied him; from the other side, the equally
unrenounceable faith in God as a person, in his guidance of the Jewish people and in the divinity of the Mosaic law.
The activity of spirit thought to be called forth by this was then an endeavour to unite these two opposed elements
into one system, and to this end the double theory is supposed to have been formed: of the allegorical interpretation
of scripture and of the
% --- p. 88 ---
\opage{88}divine intermediate beings. As one who is supposed to be borne by this conscious striving to let both these
two opposed principles hold each other in check, so that the main propositions of both remained unshaken, Philo's
doctrine is also characterized as a rational supranaturalism.

When we now in each and every respect distance ourselves from this view of God, it is not chiefly because it would
run up against a single point, but because it contradicts the whole character of the Philonic cognition. For if it has
arisen through such combining reflection, if its single representations and statements are the expression of such a
reflection, then it must also bear that reflection's marks. But that is what it does not do. By a double way this
reflection could now proceed, whereby Philo's works would indeed have acquired a very different stamp, but would yet
always have borne the character of reflection. One would rather have thought that, if Philo had really had the
difficulty before his eyes which came forward with the opposed demands touched on above, but had nevertheless regarded
it as one that vanished on closer examination --- for had it seemed to him insoluble, he would surely have contested
such a fusion and not attempted it himself --- his endeavour would then have had to aim chiefly at showing, by
reasoning comparison, how this contradiction was only apparent and was sublated --- we know, of course, by what ---
by the allegory and by the representation of the divine intermediate beings; and his philosophy would in this way
have taken a \emph{formal-logical} direction. \emph{But of all this there is in Philo no trace whatever.} As regards
\emph{the allegory}, which is supposed to be the one middle term, one will surely \ldots\ be able to infer why we can
in no way share this opinion, but nevertheless we would still direct one remark against it. If the allegory is to
have been formed for this purpose, one is surely forced to suppose that Philo, with his eye fixed on the difficulty
mentioned, with the wish to solve it if it were possible, formed this theory, not as a means to
% --- p. 89 ---
\opage{89}dazzle fools, but as that which satisfied himself, in that it did not seem to him forced by an arbitrary
power-language, but true in itself; for only then could it in turn prove the truth of the purpose in which it was
used.'' But that is not how matters stand.

That it is not ``rational supranaturalism'' but ``theosophy'' that is in question here can be seen from the
relation between the word and God.
% p. 89 from here: the setting of GL II p. 89; English from texts/nielsen/grundideernes-logik-2/translation.tex

1) \emph{The Word and God in Ambiguous Unity}. The Philonic principal category is, as is well known, the category
Logos; but whether Philo for the rest apprehends Logos as identical with God or as a being different from God, whether
this being is to be assumed to be personal or impersonal, about this the interpreters have hitherto not been able to
agree\footnote[1]{Cf.\ G.\ Steenberg. Om Synspunktet for Opfattelsen af Philos Gudserkjendelse. Kjøbenhavn 1849.
p.~84--125.}. The theosophical reflection chooses in all naivety different points of departure, without heeding the
contradictions that arise when heterogeneous ways of viewing are to be resolved into one another. On the one side
Logos is then presented as a being identical with God, as God's thinking (λογισμὸς τοῦ θεοῦ), as God's word (ῥῆμα),
God's oath (ὅρκος), God's law (νόμος), God's covenant (διαθήκη). But just as immediately Logos is set up on the other
side as a being for itself, as the bond (δεσμὸς) that binds together all things (περιέχων τὰ πάντα), which, ``in
order to hinder everything from going astray, is thought of as stretched out from the one end of the world to the
other''; Logos is ``the divine wisdom that streams through the world; the world is this wisdom's garment; wisdom
itself is ὁ τῆς φύσεως ὀρθὸς λόγος.'' Although Logos in both these apprehensions is impersonal, yet other points of
view again present themselves, from which Logos is expressly seen as a personal being, as a god subordinated to God
himself (δεύτερος θεὸς, πρωτόγονος υἱὸς), etc.

\editorial{pp.~90--262 (and p.~89 from »1) The Word and God«) are \emph{The Logic of the Fundamental Ideas II}
(1866), the same pages in the same setting: see that translation.  p.~263 below is GL II p.~263 with one line taken
up from GL II p.~264.}

% --- p. 263 ---
% English from texts/nielsen/grundideernes-logik-2/translation.tex (the same type)
\opage{263}single propositions cannot come forward with any special apodictic independence, since every relative
determination is a moment, a glimmer of the Absolute. But precisely this reflected form of the apodictic moments, a
form transparent in reflection, makes the idea, the absolutely apodictic, visible.

The general principal forms of modality: possibility, actuality and necessity, recommend themselves by turns to sound
sense in such a way that now finitude, now infinity, now the side of actuality, now the side of the idea is turned
toward the observer. At the stage of finitude it holds that there must necessarily be something that is merely
\emph{possible}, something that is \emph{actual}, and something that is \emph{necessary}. Seen under the standpoint
of infinity, the difference between possible, actual and necessary steps so far into the background that the one of
the forms involuntarily gains power over the two others, and now
\contnote{have nevertheless never been able to get the concept that is inconceivable to their thought (the concept of
infinity) separated in a satisfying manner from the method that rests upon this concept, is shown by the history of
mathematics from Leibnitz and Newton (Cf.\ Mathem.\ og Dialekt.\ p.~90--91, p.~95 ff.) down to the present time.
``In the first place, when the principle is unknown, we cannot indeed know whether all the consequences that follow
from it may likewise be admitted: we go as though blindfold, and it might well happen that after laborious exertion
we should at last stumble upon something which displayed precisely the nature of the principle, i.e.\ upon something
we could not comprehend. In the second place: is it not a shameful degradation of reason to let it follow a
connection of links, a series whose beginning and end it cannot see through?'' These scruples of Savrien's
sufficiently explain the unrest and tension in which the analysis of the infinite has constantly held mathematical
reflection (Mathem.\ Dial.\ p.~55--56), whose problem --- as a prize question set by the Berlin Academy in the year
1784 characteristically enough expresses it, ``wie aus einer widersprechenden Voraussetzung so viele wahren Sätze
haben können hergeleitet werden'' --- will not cease to return. Forelæsninger over ``Phil.\ Propæd.'' fra
Universitetsaaret 1860--61. p.~47--48.}% the Danish prints „(Uende-/hedsbegrebet)“ (misprint): translated as meant
% --- p. 264 ---
\opage{264}it is said: now that everything is possible, now that everything is actual, now that everything is
necessary. If the emphasis falls upon becoming (Heraclitus), then everything is possible; if the emphasis falls upon
being and that which is, then everything is actual (the Megarians); if the emphasis falls upon the essence of
actuality reflected into itself, the infinite substance, then everything is necessary (Spinoza). In the logical sense
necessity is overreaching, but the overreaching necessity is essentially subjective.

\begin{center}\rule{4em}{0.4pt}\\[0.5em]
C.\\[0.3em]
\textbf{Propaedeutics and the Doctrine of the Idea.}\\[0.5em]
\rule{4em}{0.4pt}\end{center}

When the idea is to be known as the idea of actuality and to take up the moments of fact into its essential content,
the development will be conditioned chiefly by a penetrating analysis not merely of the unity of the objective with
the subjective, but of the infinite richness of the determinations of objectivity themselves. As long as the
objective is merely treated in general, its peculiar character must necessarily be volatilized; for the general is
not, as has all too one-sidedly been emphasized in opposition to the psychologically subjective, the objective, but on
the contrary what is in its essence originally subjective: to treat the objective in general is to treat it in a
subjective manner, is, in other words, to treat it not as something objective but as something subjective. If, e.g.,
idea and actuality are set against each other, then the idea, although it is itself the unity of the subjective and
the objective, nevertheless immediately passes over to the subjective side, while actuality remains standing as the
particular expression of the objective; the reason for this is easy to see. The distinguishing mark of the objective
is everywhere otherness, singularity in its immediacy, the externally manifold, the diverse, etc. Now there can surely
be said in general

\editorial{pp.~265--456 are \emph{The Logic of the Fundamental Ideas I} (1864), the same pages in the same setting:
see that translation.  End of the volume.}

\end{document}
